\chapter{Prediction and Recommendations}
\label{chapter:prediction-recommendations}

\section{Regression Analysis}
\label{sec:regression}

%% ----------------------------------------------------------------------
%%  MAX 1.5 PAGES
%%  Task: make educated guesses to predict outcomes and future trends or
%%  alternative outcomes of municipal actions; validate the hypothesis;
%%  confidence intervals; graphs allowed (challenge task 7: extrapolate
%%  towards 2050, e.g. how production rate translates into backlog and
%%  risk exposure over time).
%% ----------------------------------------------------------------------

The pace of the walmuur programme determines how long the 17 walls of
Section~\ref{sec:prob} keep their timber piles, and therefore how much risk
the municipality carries towards 2050. The programme plan assumes about
100\,m of wall per year \parencite{utrecht_herstelwerven}. The team tested
this against the renewal year of each of the 32 renewed walls
(Section~\ref{sec:data}), which gives the renewed length per year in
Table~\ref{tab:production}, 3{,}431\,m in total against the 3{,}356\,m that
the municipality reports \parencite{gu2025up}.

\begin{table}[htbp]
\centering
\small
\caption{Renewed length of walmuur per year, from the renewal year of each
renewed wall and its length in the asset register.}
\label{tab:production}
\setlength{\tabcolsep}{4.5pt}
\begin{tabular}{@{}l *{12}{r}@{}}
\toprule
Year & 2015 & 2016 & 2017 & 2018 & 2019 & 2020 & 2021 & 2022 & 2023 & 2024 & 2025 & 2026 \\
Renewed [m] & 134 & 909 & 467 & 918 & 340 & 163 & 116 & 154 & 0 & 0 & 0 & 230 \\
\bottomrule
\end{tabular}
\end{table}

The first hypothesis is that the pace is constant. A linear regression
\begin{equation}
L_y = a + b\,(y - 2015) + \varepsilon_y \label{eq:production}
\end{equation}
where $L_y$ [m] is the length renewed in year $y$, $a$ [m] the intercept,
$b$ [m/yr per year] the trend and $\varepsilon_y$ [m] an independent normal
residual, gives $b = -55$\,m/yr per year (95\,\% confidence interval
\textminus105 to \textminus4, $p = 0.04$, $R^2 = 0.37$). The hypothesis of a
constant pace is therefore rejected at the 5\,\% level: the programme has
slowed down. The residuals, however, show a break rather than a gradual
decline. From 2015 to 2018 the programme renewed 607\,m per year, whereas since 2019 the pace has been 125\,m per year (95\,\%
confidence interval 22 to 228\,m), with three years without any completed
wall. Extrapolating the straight line would predict negative production
after 2027, so the line is used only as a test. The second hypothesis, that
the recent pace equals the 100\,m per year of the plan, is not rejected
($p = 0.58$), but the interval is too wide to confirm it.

To predict the backlog, the team resampled the yearly production of 2019 to
2026 (bootstrap, 20{,}000 runs) and renewed the 17 modelled walls one after
the other from 2027. The risk that a wall carries until its renewal year $T_i$
follows from its best-estimate failure probability $P_{f,i}(t)$ over the
15-year reference period (Section~\ref{sec:prob}), taken linear between 2026
and 2041:
\begin{equation}
E[N] = \sum_{i} \left[1 - \exp\!\left(-\sum_{t=2026}^{T_i} h_i(t)\right)\right],
\qquad h_i(t) = -\frac{\ln\bigl(1 - P_{f,i}(t)\bigr)}{15}, \label{eq:exposure}
\end{equation}
where $E[N]$ [--] is the expected number of walls that reach a limit state
before they are renewed and $h_i(t)$ [1/yr] the yearly hazard of wall $i$,
which assumes that failures occur at a constant rate within a reference
period. Table~\ref{tab:backlog} compares three production paces and two
orders of work: the walls with the lowest $\beta$ first, or a random order
that ignores risk.

\begin{table}[htbp]
\centering
\small
\caption{Year in which the 17 modelled walls (611\,m) are all renewed, and
expected number of walls that reach a limit state before renewal: median or
mean, with 90\,\% interval over the bootstrap.}
\label{tab:backlog}
\begin{tabular}{@{}l l r r@{}}
\toprule
Production per year & Order & All renewed & $E[N]$ \\
\midrule
Recent pace, 2019--2026 & Lowest $\beta$ first & 2031 (2029--2036) & 0.97 (0.61--1.55) \\
 & Random & 2031 (2029--2036) & 1.08 (0.66--1.73) \\
Plan, 100\,m & Lowest $\beta$ first & 2033 & 1.04 \\
 & Random & 2033 & 1.14 (1.02--1.26) \\
Fast pace, 2015--2018 & Lowest $\beta$ first & 2028 (2027--2029) & 0.56 (0.49--0.78) \\
 & Random & 2027 (2027--2029) & 0.56 (0.49--0.80) \\
\bottomrule
\end{tabular}
\end{table}

Table~\ref{tab:backlog} shows that at the recent pace the modelled walls are
renewed around 2031, well before 2050, and that about one of the 17 walls is
expected to reach a limit state before its renewal. The pace matters more
than the order: renewing the walls with the lowest $\beta$ first lowers $E[N]$
by about 10\,\%, whereas a return to the pace of 2016 to 2018 would lower it by
about 45\,\%. For the 780\,m of walmuur that the municipality still has to
restore \parencite{gu2025up}, the recent mean pace gives completion around
2033, but the lower bound of its interval, 22\,m per year, would push
completion beyond 2050. These predictions rest on few data: twelve years, of
which three without a completed wall, and a completion year that ignores
work in progress, so the bootstrap treats years as independent although
projects span several years. Moreover, $E[N]$ inherits the conservative bias of
the failure model (Section~\ref{sec:prob-validity}), and a limit state calls
for an emergency measure rather than meaning collapse. The comparison between
the scenarios is therefore more reliable than the absolute numbers;
Section~\ref{sec:results-recommendations} uses it to rank the options of the
municipality.

\clearpage
\section{Results Interpretation and Proposed Recommendations}
\label{sec:results-recommendations}

```latex
\begin{table}[htbp]
\centering
\caption{Splitting the orange class.}
\label{tab:orange_split}
\small
\begin{tabularx}{\textwidth}{>{\raggedright\arraybackslash}p{2.8cm} >{\raggedright\arraybackslash}p{3.0cm} >{\raggedright\arraybackslash}X >{\raggedright\arraybackslash}p{1.8cm}}
\toprule
\textbf{ARK critique} & \textbf{What we already have} & \textbf{What to do} & \textbf{Section} \\
\midrule
Orange mixes ``unknown'' and ``moderate risk'' &
Section~1.3 already lists $P(\text{red})$ as an output &
Split orange into \textit{orange-unknown} (non-masonry, not calculated) and \textit{orange-calculated} ($P_f$ below the limit). Use $P(\text{red})$ as a second axis: a reach with $P(\text{red}) \approx 0.5$ is where inspection yields the most. &
1.4 (flowchart), 3.2 \\
\bottomrule
\end{tabularx}
\end{table}



\begin{comment}
    The code supports three findings: the register's design life is a weak guide to which walls need work, the simple overturning model cannot separate walls, and measuring H and foundation status is worth more than refining the models.
• The register overstates the backlog. 11 of the 29 walmuren past design life are already renewed on steel (1,267 of 1,885 m), yet screening on Aanlegjaar puts them back in the queue. A foundation-status field per wall is the first fix; the Wervenkadaster would settle the 8 unclear walls.
• The block model is too conservative to rank walls. Mean load 2.78–4.22 kNm/m against mean resistance 2.09 kNm/m: it fails walls that stand. This fits the brief's point that code-based models understate capacity. The missing pieces are the floor, the soil on a heel and the piles.
• The pile model supports relative statements only. Its β of 3.24 at 15 years rests on Amsterdam inputs, and it does not reproduce its own benchmark (β 1.5).
• The safety level is a choice that changes the answer. The same β of 3.24 passes CC2 and fails CC3.
• What to measure next, in order of effect on the model:
    1. Canal bottom per rak (soundings): sets H, and the soil load grows with H³.
    2. The load case on the werf: where people, terraces and vehicles stand (a).
    3. K near the walls: borings within 150 m agree for only 9 of the 51 walls.
    4. Wall-thickness profile and pile diameter: the resistance side.
• Candidates for the one data request. The code names three: canal-bottom soundings per rak, the foundation renewal year per rak, and the Wervenkadaster.
• Not in the code: costs, a programming direction, and synergies with other works.
\end{comment}
```

% This needs `\usepackage{booktabs, tabularx, array, amsmath}` in your preamble.
%% ----------------------------------------------------------------------
%%  MAX 1.5 PAGES
%%
%%  TASK DESCRIPTION
%%  Discuss the obtained solution by linking the results to the level of
%%  abstraction of the selected models, the data uncertainty, and the
%%  assumptions made in Section \ref{sec:assumptions-data}. Discuss and
%%  justify the soundness of the results.
%%  You may include graphs to support statements, and use the appendices
%%  for other supporting information.
%%  Based on the interpretation of the results, what is the recommendation
%%  you give to the Municipality of Utrecht? Communicate the results to
%%  non-experts in data and uncertainty management. You may suggest
%%  feasible interventions, expected benefits, associated costs,
%%  uncertainties and risks, as well as potential synergies with other
%%  transport and spatial planning initiatives.
%%
%%  CHECKLIST (challenge task 8: prepare actionable recommendations)
%%   - Integrate the analyses into a substantiated recommendation stating
%%     what is known, what is assumed, which programming direction you
%%     propose, and what the municipality should measure next in order to
%%     reduce the uncertainty that currently dominates the decision.
%%   - Flag which recommendations need an institutional or regulatory
%%     change rather than money.
%% ----------------------------------------------------------------------

